\documentclass[12pt, twoside]{article}
\input{../previous-materials/preamble}
\usepackage{float}

\title{Physics}
\author{Chris Huson}
\date{May 2026}

\fancyhead[LE]{\thepage}
\fancyhead[RO]{\thepage \\ First \& last name: \hspace{2.25cm} \,\\ Grade: \hspace{2.25cm} \,}
\fancyhead[LO]{La Scuola d'Italia / Huson / Physics \\* 11 May 2026}

\begin{document}

\subsubsection*{Relativity Follow-Up Worksheet}

\textbf{Useful constants and units:} \\
$c = 3.00\times10^{8}~\text{m/s}$ \quad Earth orbital speed $\approx 30~\text{km/s}$ \quad $g = 9.8~\text{m/s}^{2}$ \\
$1~\text{km} = 1000~\text{m}$ \quad $1~\text{min} = 60~\text{s}$ \quad $1~\text{h} = 3600~\text{s}$ \quad $1~\text{day} = 86{,}400~\text{s}$ \\
$1~\text{J} = 1~\text{N}\cdot\text{m} = 1~\text{kg}\cdot\text{m}^{2}/\text{s}^{2}$ \quad $1~\text{megaton TNT} \approx 4.184\times10^{15}~\text{J}$ \\
$1~\text{eV} = 1.6\times10^{-19}~\text{J}$ \quad $1~\text{MeV} = 10^{6}~\text{eV}$ \quad Proton rest mass $\approx 1.67\times10^{-27}~\text{kg}$

\vspace{0.3cm}

\begin{enumerate}[itemsep=0.5cm]

\item \textbf{Reference frames.} In your own words, explain what a ``frame of reference'' is. Give one everyday example where two observers measure different values for the same motion (for example, a passenger on a train and a person on the platform).
  \vspace{5cm}

\item \textbf{Michelson--Morley.} Briefly describe what Michelson and Morley were trying to detect, and what they actually found.
  \begin{enumerate}[itemsep=2cm]
    \item What was the ``ether'' supposed to be?
    \item Why did scientists expect the speed of light to depend on Earth's motion through it?
    \item What did the null result of the experiment imply about the speed of light?
  \end{enumerate}

\newpage
\item \textbf{Light speed in every frame.} A spaceship moving at half the speed of light shines a flashlight forward. An observer on a nearby planet measures the speed of the light from the flashlight. What value does the planet observer measure, and why?
  \vspace{6cm}

\item \textbf{Unit conversions with large speeds.}
  \begin{enumerate}[itemsep=2.25cm]
    \item Convert Earth's orbital speed, $30~\text{km/s}$, into m/s. Write the answer in scientific notation.
    \item Convert the speed of light, $3.00\times10^{8}~\text{m/s}$, into km/s.
    \item How many times faster than Earth's orbital speed is the speed of light? Show the ratio $c / (30~\text{km/s})$.
  \end{enumerate}

\newpage
\item \textbf{How far does Earth travel?} Using $30~\text{km/s}$ as Earth's orbital speed:
  \begin{enumerate}[itemsep=3cm]
    \item How far does Earth move along its orbit in one hour? (Give the answer in km.)
    \item How far does Earth move along its orbit in one day? (Give the answer in km, in scientific notation.)
  \end{enumerate}
  \vspace{0.5cm}

\item \textbf{How far does light travel?} Using $c = 3.00\times10^{8}~\text{m/s}$:
  \begin{enumerate}[itemsep=2cm]
    \item How far does light travel in $1.00~\text{s}$? (Give the answer in m and in km.)
    \item How far does light travel in $1.00~\text{min}$?
    \item Estimate how far light travels in one year (a ``light-year'') in meters. Use $1~\text{year} \approx 3.15\times10^{7}~\text{s}$.
  \end{enumerate}

\newpage
\item \textbf{Work and energy: climbing to class.} The \emph{joule} (J) is the SI unit of energy and of work. The work done lifting an object of mass $m$ straight up by a height $h$ against gravity is
$W = mgh$, which equals the gain in gravitational potential energy.

  \begin{enumerate}[itemsep=1.8cm]
    \item Estimate your own mass in kg. Using $g = 9.8~\text{m/s}^{2}$, compute the work, in joules, required to lift yourself $h = 20~\text{m}$ (roughly the climb from the street up to our sixth-floor classroom).
    \item A typical chocolate bar stores about $1.0\times10^{6}~\text{J}$ of food energy. How many trips up to the classroom could that energy power, in principle?
    \item Show that $1~\text{kg}\cdot\text{m}^{2}/\text{s}^{2}$ is the same combination of units as $1~\text{N}\cdot\text{m}$. (Hint: $\text{N} = \text{kg}\cdot\text{m}/\text{s}^{2}$.)
  \end{enumerate}
  \vspace{0.3cm}

\item \textbf{Mass--energy equivalence.} Einstein's relation is $E = mc^{2}$.
  \begin{enumerate}[itemsep=1.8cm]
    \item Compute the energy released if $m = 1.00~\text{g} = 1.00\times10^{-3}~\text{kg}$ is converted entirely into energy. Give the answer in joules, in scientific notation.
    \item Convert that energy into megatons of TNT. (Use $1~\text{megaton} \approx 4.184\times10^{15}~\text{J}$.)
    \item In one sentence, explain why such a small mass releases such a large amount of energy.
  \end{enumerate}

\newpage
\item \textbf{Length contraction.} A spaceship measured at rest is $100~\text{m}$ long. It now flies past Earth at a speed close to $c$.
  \begin{enumerate}[itemsep=2cm]
    \item According to observers on Earth, is the moving spaceship measured to be longer than, shorter than, or the same length as $100~\text{m}$? Explain.
    \item According to the astronauts inside the ship, what length do \emph{they} measure for their own ship? Explain.
    \item In which direction does the contraction occur --- along the direction of motion, perpendicular to it, or both?
  \end{enumerate}
  \vspace{0.4cm}

\item \textbf{Time dilation.} A pair of identical clocks start out synchronized. One stays on Earth; the other is flown around the world in a fast jet and then returned.
  \begin{enumerate}[itemsep=1.8cm]
    \item Which clock, if either, will read \emph{less} elapsed time when they are compared? Why?
    \item Why don't we notice time dilation when we ride in a car or an airplane in everyday life?
    \item Muons created high in the atmosphere by cosmic rays are unstable and would normally decay long before reaching the ground, yet many reach sea level. How does time dilation help explain this?
  \end{enumerate}

\newpage
\item \textbf{From mass to MeV.} Physicists often measure tiny energies in \emph{electron-volts}. One eV is the energy an electron gains falling through one volt: $1~\text{eV} = 1.6\times10^{-19}~\text{J}$, and $1~\text{MeV} = 10^{6}~\text{eV} = 1.6\times10^{-13}~\text{J}$.

  \begin{enumerate}[itemsep=1.2cm]
    \item Using $E = mc^{2}$ with the proton's rest mass $m_p = 1.67\times10^{-27}~\text{kg}$, compute the proton's rest energy in joules. Give your answer in scientific notation.
    \item Convert that energy into MeV. Your answer should be close to $940~\text{MeV}$.
    \item In your own words, what does it mean to say ``a proton has a rest energy of about $938~\text{MeV}$''?
  \end{enumerate}
  \vspace{0.2cm}

\item \textbf{Particle accelerators (the LHC).} The Large Hadron Collider is a circular particle accelerator at CERN, near Geneva on the French--Swiss border. Its main ring is about $27~\text{km}$ around, sits roughly $100~\text{m}$ underground, was built between $1998$ and $2008$ at a cost of around \$5 billion, and is operated by CERN (the European Organization for Nuclear Research, funded by more than $20$ member countries). The Higgs boson was announced there in $2012$. At the LHC, protons reach a total energy about $7000$ times their rest energy.

  \begin{enumerate}[itemsep=1cm]
    \item If a proton's rest energy is about $938~\text{MeV}$, estimate its total energy at the LHC. Give your answer in MeV and in TeV ($1~\text{TeV} = 10^{6}~\text{MeV}$).
    \item Explain in one or two sentences why it becomes harder and harder to speed up a proton as it gets closer to the speed of light.
    \item Can a proton (or any object with rest mass) ever actually reach the speed of light? Why or why not?
  \end{enumerate}
  \vspace{0.2cm}

\item \textbf{Discussion / short answer.} Choose \emph{one} of the following and answer in 3--5 sentences.
  \begin{enumerate}[itemsep=0pt, parsep=0pt, topsep=0pt]
    \item Why is the speed of light considered a universal speed limit, and what would the consequences be if it weren't?
    \item If you traveled in a spaceship at close to the speed of light for what felt like one year, more than one year would pass on Earth. Would you consider this a form of time travel? Explain.
    \item How does the Michelson--Morley experiment connect to Einstein's claim that the laws of physics look the same in every inertial frame?
  \end{enumerate}

\end{enumerate}

\end{document}
